\chapter{特征值与特征向量}

本章研究方阵最内在的结构：哪些方向在变换下\textbf{保持不变}？答案是特征向量；
不变的倍数是特征值。作为准备，先介绍内积、长度与正交性；随后讨论相似矩阵
与对角化；最后证明实对称矩阵必可正交对角化——它是第六章化简二次型的钥匙。

\section{向量的内积、长度与正交性}

\begin{definition}\label{def:inner-product}
设 $\alpha=(a_1,\cdots,a_n)$、$\beta=(b_1,\cdots,b_n)\in\R^n$，称数
\[
  (\alpha,\beta)=a_1b_1+a_2b_2+\cdots+a_nb_n=\alpha\trans\beta
\]
为 $\alpha$ 与 $\beta$ 的\textbf{内积}。称 $\|\alpha\|=\sqrt{(\alpha,\alpha)}
=\sqrt{a_1^2+\cdots+a_n^2}$ 为 $\alpha$ 的\textbf{长度}（范数）；长度为 $1$ 的
向量称为\textbf{单位向量}。非零向量 $\alpha$ 的单位化向量为
$\alpha_0=\alpha/\|\alpha\|$。
\end{definition}

内积满足对称性 $(\alpha,\beta)=(\beta,\alpha)$、对两个变元的线性性，以及
$(\alpha,\alpha)\ge0$ 且等号成立当且仅当 $\alpha=0$。

\begin{theorem}[柯西--施瓦茨不等式]\label{thm:schwarz}
$|(\alpha,\beta)|\le\|\alpha\|\,\|\beta\|$，等号成立当且仅当 $\alpha,\beta$
线性相关。
\end{theorem}
\begin{proof}
若 $\beta=0$ 显然。设 $\beta\neq0$，对任意实数 $t$，
$0\le(\alpha-t\beta,\alpha-t\beta)=\|\alpha\|^2-2t(\alpha,\beta)+t^2\|\beta\|^2$。
右端是 $t$ 的二次函数且恒非负，故判别式 $\le0$：
$(\alpha,\beta)^2\le\|\alpha\|^2\|\beta\|^2$。
\end{proof}

由柯西--施瓦茨不等式，可以定义两向量夹角 $\theta$ 的余弦
$\cos\theta=(\alpha,\beta)/(\|\alpha\|\|\beta\|)$。当 $(\alpha,\beta)=0$ 时
称 $\alpha$ 与 $\beta$ \textbf{正交}（垂直）。一组两两正交的非零向量称为
\textbf{正交向量组}；若其中每个还是单位向量，称为\textbf{标准正交组}。

\begin{theorem}\label{thm:orthogonal-independent}
正交向量组必线性无关。
\end{theorem}
\begin{proof}
设 $\beta_1,\cdots,\beta_m$ 两两正交且 $k_1\beta_1+\cdots+k_m\beta_m=0$。两边
与 $\beta_j$ 作内积，利用两两正交性得 $k_j\|\beta_j\|^2=0$；$\beta_j\neq0$，
故 $k_j=0$（$j=1,\cdots,m$）。
\end{proof}

\begin{qa}{怎样把一组无关向量“改造”成正交组？减掉的到底是什么？}
答：用\textbf{施密特正交化}。设 $\alpha_1,\cdots,\alpha_m$ 线性无关，令
\[
  \beta_1=\alpha_1,\qquad
  \beta_j=\alpha_j-\sum_{i<j}\frac{(\alpha_j,\beta_i)}{\|\beta_i\|^2}\,\beta_i
  \quad(j=2,\cdots,m).
\]
减掉的第 $i$ 项正是 $\alpha_j$ 在已有正交方向 $\beta_i$ 上的\textbf{投影}
（图~\ref{fig:schmidt}）：$\dfrac{(\alpha_j,\beta_i)}{\|\beta_i\|^2}\beta_i$
是 $\beta_i$ 方向上“与 $\alpha_j$ 内积相同”的那个向量。减完全部投影后，剩余
分量与每个 $\beta_i$ 的内积都是 $0$——垂直是“减掉投影”的直接后果：
\[
  (\beta_j,\beta_i)
  =(\alpha_j,\beta_i)-\frac{(\alpha_j,\beta_i)}{\|\beta_i\|^2}(\beta_i,\beta_i)=0
  \quad(i<j).
\]
最后逐个单位化，即得与原向量组\textbf{张成相同}的标准正交组（每一步只做减法
与数乘，不改变张成空间，故“改造”是无损的）。
\end{qa}

\begin{figure}[htbp]
\centering
\begin{tikzpicture}[>=Stealth, scale=1.15]
  \draw[->, gray!70] (-0.9,0) -- (2.6,0);
  \draw[->, gray!70] (0,-0.5) -- (0,2.3);
  \draw[very thick, blue!60!black, ->] (0,0) -- (1.8,0.6)
        node[below right] {$\beta_1=\alpha_1$};
  \draw[thick, gray!70!black, ->] (0,0) -- (1.2,1.8) node[above] {$\alpha_2$};
  \draw[dashed, blue!60!black, ->] (0,0) -- (1.62,0.54)
        node[below, xshift=2pt] {\small 投影};
  \draw[dashed, gray] (1.2,1.8) -- (1.62,0.54);
  \draw[very thick, red!70!black, ->] (0,0) -- (-0.42,1.26)
        node[left] {$\beta_2=\alpha_2-$ 投影};
  \draw[dashed, red!70!black] (1.62,0.54) -- (1.2,1.8);
  \draw[gray!70] (1.62,0.54) ++(18:0.22) arc[start angle=18, end angle=90, radius=0.22];
\end{tikzpicture}
\caption{施密特正交化：从 $\alpha_2$ 中减去它在 $\beta_1$ 上的投影，剩余分量
$\beta_2$ 与 $\beta_1$ 垂直（图中直角标记处）。}
\label{fig:schmidt}
\end{figure}

\begin{example}\label{ex:schmidt}
把 $\alpha_1=(1,1,0)$、$\alpha_2=(1,0,1)$、$\alpha_3=(0,1,1)$ 正交化并
单位化。
\end{example}
\begin{proof}[解]
\[
  \beta_1=(1,1,0);
  \qquad
  \beta_2=\alpha_2-\tfrac12\beta_1=(\tfrac12,-\tfrac12,1);
  \]
  \[
  \beta_3=\alpha_3-\tfrac12\beta_1-\tfrac13\beta_2
         =\alpha_3-\tfrac12(1,1,0)-\tfrac13(\tfrac12,-\tfrac12,1)
         =(-\tfrac23,\tfrac23,\tfrac23).
\]
去掉正倍数，取 $\beta_1'=(1,1,0)$，$\beta_2'=(1,-1,2)$，$\beta_3'=(-1,1,1)$，
两两正交（可逐一验证内积为零），单位化得
\[
  \gamma_1=\tfrac{1}{\sqrt2}(1,1,0),\quad
  \gamma_2=\tfrac{1}{\sqrt6}(1,-1,2),\quad
  \gamma_3=\tfrac{1}{\sqrt3}(-1,1,1).
\]
\end{proof}

\begin{definition}\label{def:orthogonal-matrix}
满足 $A\trans A=E$ 的 $n$ 阶方阵 $A$ 称为\textbf{正交矩阵}。
\end{definition}

\begin{theorem}\label{thm:orthogonal-matrix}
$A$ 为正交矩阵 $\iff$ $A$ 的列向量组是标准正交组。正交矩阵可逆且
$A^{-1}=A\trans$，$|A|=\pm1$；两个正交矩阵之积仍正交。正交变换保持内积、
长度与夹角不变。
\end{theorem}
\begin{proof}
$A\trans A=E$ 按元素写开即“第 $i$ 列与第 $j$ 列的内积为 $\delta_{ij}$”。
取行列式得 $|A|^2=1$。保持内积：
$(Ax,Ay)=(Ax)\trans(Ay)=x\trans A\trans A\,y=(x,y)$，长度与夹角由内积决定。
\end{proof}

\section{特征值与特征向量}\label{sec:eigen}

\begin{definition}\label{def:eigen}
设 $A$ 为 $n$ 阶方阵。若存在数 $\lambda$ 与\textbf{非零}向量 $\xi$，使
\[
  A\xi=\lambda\xi,
\]
则称 $\lambda$ 是 $A$ 的一个\textbf{特征值}，$\xi$ 是属于 $\lambda$ 的
\textbf{特征向量}。
\end{definition}

\begin{insight}[本质：特征向量 = 变换“不动”的方向]
把 $A$ 看作空间的变换：一般的向量被移动后\textbf{连方向都变了}；特征向量是
幸运者——它只被伸缩，不偏转，特征值 $\lambda$ 就是伸缩的倍数（$\lambda<0$
表示反向）。于是“找特征向量”就是问：\textbf{这个变换有哪些天生不动的方向，
各伸缩多少倍？}（图~\ref{fig:eig-geo}）。若能找到 $n$ 个无关的不动方向，变换
在这组“主轴”下就只是一组独立的伸缩——这正是“对角化”的全部含义。
\end{insight}

\begin{figure}[htbp]
\centering
\begin{tikzpicture}[>=Stealth, scale=0.9]
  \begin{scope}
    \draw[->, gray!70] (-0.4,0) -- (2.6,0);
    \draw[->, gray!70] (0,-0.4) -- (0,2.6);
    \draw[thick, gray!70!black, ->] (0,0) -- (0.8,1.4) node[above left] {$v$};
    \draw[very thick, red!70!black, ->] (0,0) -- (1.6,1.4) node[right] {$Av$};
    \node[below] at (1.1,-0.7) {一般向量：方向被改变};
  \end{scope}
  \draw[->, very thick, gray] (3.6,1.3) -- (4.9,1.3)
        node[midway, above] {$A=\operatorname{diag}(2,1)$};
  \begin{scope}[xshift=6.0cm]
    \draw[->, gray!70] (-0.4,0) -- (2.9,0);
    \draw[->, gray!70] (0,-0.4) -- (0,2.6);
    \draw[thick, blue!60!black, ->] (0,0) -- (1.3,0);
    \draw[very thick, blue!60!black, ->] (0,0) -- (2.6,0)
        node[below] {$\xi$（$\lambda=2$）};
    \draw[very thick, red!70!black, ->] (0,0) -- (0,1.6)
        node[left] {$\eta$（$\lambda=1$）};
    \node[below] at (1.2,-0.8) {特征向量：只伸缩，不偏转};
  \end{scope}
\end{tikzpicture}
\caption{伸缩变换 $\operatorname{diag}(2,1)$：一般向量（左）方向改变；坐标轴
方向（右）天生不动，$x$ 轴方向拉伸 $2$ 倍，$y$ 轴方向保持 $1$ 倍。}
\label{fig:eig-geo}
\end{figure}

\begin{theorem}\label{thm:eigen-criterion}
$\lambda$ 是 $A$ 的特征值 $\iff$ $|A-\lambda E|=0$；$\xi$ 是属于 $\lambda$ 的
特征向量 $\iff$ $\xi$ 是齐次方程组 $(A-\lambda E)x=0$ 的非零解。
\end{theorem}
\begin{proof}
$A\xi=\lambda\xi$ 即 $(A-\lambda E)\xi=0$。它有非零解当且仅当系数矩阵降秩，
即 $|A-\lambda E|=0$（第四章定理~\ref{thm:fundamental} 的特例）。
\end{proof}

称 $f(\lambda)=|A-\lambda E|$ 为 $A$ 的\textbf{特征多项式}。求特征值的步骤：
解特征方程 $f(\lambda)=0$；对每个特征值解 $(A-\lambda E)x=0$，其基础解系给
出全部特征向量。

\begin{example}\label{ex:eigen-2}
求 $A=\begin{pmatrix} 3 & -1 \\ -1 & 3 \end{pmatrix}$ 的特征值与特征向量。
\end{example}
\begin{proof}[解]
\[
  |A-\lambda E|=\begin{vmatrix} 3-\lambda & -1 \\ -1 & 3-\lambda \end{vmatrix}
  =(3-\lambda)^2-1=(2-\lambda)(4-\lambda),
\]
故 $\lambda_1=2$，$\lambda_2=4$。对 $\lambda_1=2$：解 $(A-2E)x=0$，即
$\begin{pmatrix} 1 & -1 \\ -1 & 1 \end{pmatrix}x=0$，得 $\xi_1=(1,1)\trans$
（验证：$A\xi_1=(2,2)\trans=2\xi_1$）。对 $\lambda_2=4$：得
$\xi_2=(1,-1)\trans$（$A\xi_2=(4,-4)\trans=4\xi_2$）。
\end{proof}

\begin{theorem}\label{thm:eigen-sum-product}
设 $n$ 阶方阵 $A$ 的全部特征值为 $\lambda_1,\cdots,\lambda_n$（按重数计），
则
\[
  \lambda_1+\cdots+\lambda_n=\operatorname{tr}(A)=a_{11}+\cdots+a_{nn},
  \qquad
  \lambda_1\lambda_2\cdots\lambda_n=|A|.
\]
\end{theorem}

在特征多项式的展开中，$\lambda^{n-1}$ 的系数由主对角线乘积这一项唯一贡献，
故等于 $-\operatorname{tr}(A)$；令 $\lambda=0$ 得常数项 $f(0)=|A|$；再由根与
系数的关系即得（详细证明从略）。该定理是很好的快速校验：例~\ref{ex:eigen-2}
中 $2+4=6=\operatorname{tr}(A)$，$2\times4=8=|A|$，两条都吻合。

\begin{theorem}\label{thm:distinct-independent}
属于\textbf{互不相同}特征值的特征向量线性无关。
\end{theorem}
\begin{proof}
设 $\lambda_1,\cdots,\lambda_m$ 互不相同，$\xi_i$ 属于 $\lambda_i$，且
$c_1\xi_1+\cdots+c_m\xi_m=0$。对等式两边施加乘积
$(A-\lambda_2E)(A-\lambda_3E)\cdots(A-\lambda_mE)$。因
$(A-\lambda_jE)\xi_j=0$，第 $j\ ( \ge2)$ 项全部被“杀死”；第 1 项变为
\[
  c_1\prod_{j=2}^{m}(\lambda_1-\lambda_j)\,\xi_1=0 .
\]
$\xi_1\neq0$ 且各因子非零，故 $c_1=0$。同理（依次换序施加）可得每个
$c_i=0$。
\end{proof}

\begin{qa}{特征向量为什么必须非零？$\lambda=0$ 又说明什么？}
答：零向量对\textbf{任何} $\lambda$ 都满足 $A\cdot0=\lambda\cdot0$，若允许
它当特征向量，“特征”二字就毫无信息量，故定义明确要求 $\xi\neq0$。$\lambda=0$
成为特征值 $\iff|A-0\cdot E|=|A|=0$，即 $A$ 奇异；此时特征向量正是齐次方程组
$Ax=0$ 的非零解——“被变换压成零的方向”。这是第二章“等价网”的又一条边：
\textbf{有零特征值 $\iff$ 不可逆 $\iff$ 秩亏 $\iff$ 列相关 $\iff$ 压扁}。
\end{qa}

\section{相似矩阵与矩阵的对角化}\label{sec:similar}

\begin{definition}\label{def:similar}
设 $A$、$B$ 为 $n$ 阶方阵。若存在可逆矩阵 $P$ 使 $P^{-1}AP=B$，则称 $B$ 与
$A$ \textbf{相似}，记作 $A\sim B$。若 $B$ 是对角矩阵 $\Lambda$，则称 $A$
\textbf{可对角化}。
\end{definition}

相似满足反身性、对称性、传递性；相似的矩阵同秩、同行列式、同迹、同特征
多项式（从而同特征值）。同特征多项式的验证：
\[
  |B-\lambda E|=|P^{-1}AP-\lambda E|
  =|P^{-1}(A-\lambda E)P|=|P^{-1}|\,|A-\lambda E|\,|P|=|A-\lambda E|.
\]
注意反之不成立：相似比“同特征多项式”苛刻得多。

\begin{insight}[本质：相似 = 同一个变换在不同基下的两张照片]
把 $A$ 看作一个变换，取一组新基排成矩阵 $P$，则同一个变换在新基下的“照片”
恰是 $P^{-1}AP$。相似的矩阵不是两个变换，而是\textbf{一个变换的两个机位}。
对角化就是挑一个最好的机位——让变换看起来最简单：沿 $n$ 个互相独立的方向
各自伸缩 $\lambda_i$，别无杂动作。这个“最好的机位”存在当且仅当 $A$ 有
$n$ 个线性无关的特征向量（下面的定理）。
\end{insight}

\begin{theorem}\label{thm:diagonalizable}
$n$ 阶方阵 $A$ 可对角化 $\iff$ $A$ 有 $n$ 个线性无关的特征向量。此时令
$P=(\xi_1,\cdots,\xi_n)$（列为特征向量），则
\[
  P^{-1}AP=\Lambda=\operatorname{diag}(\lambda_1,\cdots,\lambda_n),
\]
且 $\Lambda$ 的第 $i$ 个对角元属于 $P$ 的第 $i$ 列。特别地，若 $A$ 有 $n$ 个
互不相同的特征值，则 $A$ 必可对角化（逆命题不成立，如 $E$ 本身）。
\end{theorem}
\begin{proof}
把“$A\xi_i=\lambda_i\xi_i$（$i=1,\cdots,n$）”拼成矩阵等式：
$AP=P\Lambda$，即 $P^{-1}AP=\Lambda$（$P$ 的列无关保证 $P$ 可逆）。反之，
由 $P^{-1}AP=\Lambda$ 得 $AP=P\Lambda$，逐列读出 $A\xi_i=\lambda_i\xi_i$，
且 $P$ 可逆保证这 $n$ 个特征向量无关。“互不相同”情形由定理
~\ref{thm:distinct-independent} 立得。
\end{proof}

\begin{qa}{$P$ 的列必须放什么？$\Lambda$ 的对角元顺序可以随便吗？}
答：$P$ 的第 $i$ 列\textbf{必须}是属于 $\lambda_i$ 的特征向量——由证明中的
矩阵等式 $AP=P\Lambda$ 逐列对照即知。因此 $\Lambda$ 的对角元与 $P$ 的列
\textbf{成对绑定}：交换 $\Lambda$ 中两个特征值的位置，必须同时交换 $P$ 中
对应的两个列。单独换其中一边，等式立刻失效。
\end{qa}

\begin{example}\label{ex:diagonalize}
判断 $A=\begin{pmatrix} 4 & 0 & 0 \\ 0 & 3 & 1 \\ 0 & 1 & 3 \end{pmatrix}$
能否对角化；若能，求 $P$ 与 $\Lambda$，并计算 $A^k$。
\end{example}
\begin{proof}[解]
\[
  |A-\lambda E|=(4-\lambda)\bigl[(3-\lambda)^2-1\bigr]
  =(4-\lambda)^2(2-\lambda),
\]
$\lambda_1=\lambda_2=4$，$\lambda_3=2$。对 $\lambda=4$：解 $(A-4E)x=0$，
得 $x_2=x_3$，基础解系 $\xi_1=(1,0,0)\trans$，$\xi_2=(0,1,1)\trans$；对
$\lambda=2$：得 $\xi_3=(0,1,-1)\trans$。三个特征向量线性无关（其中
$\xi_1,\xi_2$ 属于 $\lambda=4$，$\xi_3$ 属于 $\lambda=2$，跨特征值由定理
~\ref{thm:distinct-independent} 保证无关，同特征值内可直接验证无关）。故
\[
  P=\begin{pmatrix} 1 & 0 & 0 \\ 0 & 1 & 1 \\ 0 & 1 & -1 \end{pmatrix},
  \qquad
  P^{-1}AP=\begin{pmatrix} 4 & & \\ & 4 & \\ & & 2 \end{pmatrix}=\Lambda .
\]
由 $A=P\Lambda P^{-1}$ 得 $A^k=P\Lambda^kP^{-1}$，且
$P^{-1}=\begin{pmatrix} 1 & 0 & 0 \\ 0 & \tfrac12 & \tfrac12 \\
0 & \tfrac12 & -\tfrac12 \end{pmatrix}$，于是
\[
  A^k=\begin{pmatrix}
    4^k & 0 & 0 \\
    0 & \dfrac{4^k+2^k}{2} & \dfrac{4^k-2^k}{2} \\
    0 & \dfrac{4^k-2^k}{2} & \dfrac{4^k+2^k}{2}
  \end{pmatrix}.
\]
（取 $k=1$ 可核对与 $A$ 一致。）
\end{proof}

\section{实对称矩阵的正交对角化}\label{sec:symmetric}

一般方阵未必可对角化，但实对称矩阵是一个例外，而且可以要求变换矩阵是
\textbf{正交}矩阵。

\begin{theorem}\label{thm:symmetric-real}
实对称矩阵的特征值都是实数。
\end{theorem}
\begin{proof}
设 $A\trans=A$（实矩阵），$A\xi=\lambda\xi$，$\xi\neq0$（暂允许复向量）。
取共轭转置：$\bar{\xi}\trans A\trans=\bar{\lambda}\bar{\xi}\trans$；由 $A$
实且对称，即 $\bar{\xi}\trans A=\bar{\lambda}\bar{\xi}\trans$。于是
\[
  \bar{\xi}\trans A\xi=\lambda\,\bar{\xi}\trans\xi,
  \qquad
  \bar{\xi}\trans A\xi=(\bar{\xi}\trans A)\xi=\bar{\lambda}\,\bar{\xi}\trans\xi,
\]
两式相减得 $(\lambda-\bar{\lambda})\,\bar{\xi}\trans\xi=0$。而
$\bar{\xi}\trans\xi=\sum|x_i|^2>0$，故 $\lambda=\bar{\lambda}$ 为实数。
\end{proof}

\begin{theorem}\label{thm:symmetric-orthogonal}
实对称矩阵属于不同特征值的特征向量彼此正交。
\end{theorem}
\begin{proof}
设 $A\xi_1=\lambda_1\xi_1$，$A\xi_2=\lambda_2\xi_2$，$\lambda_1\neq\lambda_2$。
则
\[
  \lambda_1(\xi_1,\xi_2)=\xi_2\trans A\xi_1
  =\xi_2\trans A\trans\xi_1=(A\xi_2)\trans\xi_1=\lambda_2(\xi_1,\xi_2),
\]
故 $(\lambda_1-\lambda_2)(\xi_1,\xi_2)=0$，即 $(\xi_1,\xi_2)=0$。
\end{proof}

\begin{theorem}[谱定理，证明从略]\label{thm:spectral}
设 $A$ 为 $n$ 阶实对称矩阵，则存在正交矩阵 $Q$，使
\[
  Q\trans A\,Q=Q^{-1}AQ=\Lambda=\operatorname{diag}(\lambda_1,\cdots,\lambda_n),
\]
其中 $\lambda_i$ 为 $A$ 的全部特征值（按重数计）。
\end{theorem}

操作步骤：求出全部特征值；对每个特征值求 $(A-\lambda E)x=0$ 的基础解系
（同一特征值内的解向量用施密特正交化——跨特征值的自动正交）；全部单位化
后排成 $Q$ 的列。

\begin{insight}[本质：实对称 = 沿互相垂直的主轴纯拉伸]
谱定理的几何内容：对称变换没有任何“旋转、剪切”的杂动作，它只是沿 $n$ 条
\textbf{互相垂直的主轴}各自伸缩。对比一下损失在哪：旋转矩阵
$\begin{pmatrix} 0 & -1 \\ 1 & 0 \end{pmatrix}$ 把每个方向都转走 $90^\circ$，
一个实特征方向都没有——所以“不可对角化/没有实特征值”并不可耻，只是这类
变换的“不动方向”不在这个实平面上。
\end{insight}

\begin{example}\label{ex:spectral}
设 $A=\begin{pmatrix} 0 & 1 & 1 \\ 1 & 0 & 1 \\ 1 & 1 & 0 \end{pmatrix}$，
求正交矩阵 $Q$ 使 $Q\trans AQ$ 为对角矩阵。
\end{example}
\begin{proof}[解]
$|A-\lambda E|=-(\lambda-2)(\lambda+1)^2$（各行加到第一行提出
$(2-\lambda)$ 再化简），$\lambda_1=2$，$\lambda_{2,3}=-1$（二重）。对
$\lambda=2$：解 $(A-2E)x=0$ 得 $\xi_1=(1,1,1)\trans$。对 $\lambda=-1$：
解 $(A+E)x=0$，即 $x_1+x_2+x_3=0$，取 $\eta_2=(1,-1,0)\trans$，
$\eta_3=(1,0,-1)\trans$，施密特正交化：
\[
  \beta_2=\eta_2=(1,-1,0)\trans,\qquad
  \beta_3=\eta_3-\tfrac12\beta_2=(\tfrac12,\tfrac12,-1)\trans .
\]
单位化：
\[
  Q=\left(
  \tfrac{1}{\sqrt3}\xi_1,\ \tfrac{1}{\sqrt2}\beta_2,\ \tfrac{1}{\sqrt6}(1,1,-2)\trans
  \right)
  =\begin{pmatrix}
    \tfrac{1}{\sqrt3} & \tfrac{1}{\sqrt2} & \tfrac{1}{\sqrt6} \\
    \tfrac{1}{\sqrt3} & -\tfrac{1}{\sqrt2} & \tfrac{1}{\sqrt6} \\
    \tfrac{1}{\sqrt3} & 0 & -\tfrac{2}{\sqrt6}
  \end{pmatrix},
\]
则 $Q\trans AQ=\operatorname{diag}(2,-1,-1)$（$Q$ 的列标准正交可逐一验证）。
\end{proof}

\chapterexercises{五}

\begin{enumerate}[label=\arabic*., itemsep=2pt]
\item 求 $\begin{pmatrix} 2 & 1 \\ 1 & 2 \end{pmatrix}$ 的特征值与特征向量。
\item 把 $\alpha_1=(1,1,1)$、$\alpha_2=(1,2,3)$ 施密特正交化并单位化。
\item 判断下列矩阵是否正交矩阵：
  $\text{(1)}\ \dfrac{1}{\sqrt2}\begin{pmatrix} 1 & -1 \\ 1 & 1 \end{pmatrix}$；
  \quad $\text{(2)}\ \begin{pmatrix} 1 & 1 \\ 0 & 1 \end{pmatrix}$。
\item 求 $\begin{pmatrix} 2 & 0 & 0 \\ 0 & 3 & 1 \\ 0 & 1 & 3 \end{pmatrix}$
  的特征值与特征向量，并判断它是否可对角化。
\item 利用第 1 题的结果计算 $\begin{pmatrix} 2 & 1 \\ 1 & 2 \end{pmatrix}^{10}$。
\item 求正交矩阵 $Q$ 使 $Q\trans\begin{pmatrix} 0 & 1 \\ 1 & 0 \end{pmatrix}Q$
  为对角矩阵。
\item 证明：$A$ 与 $A\trans$ 有相同的特征多项式。
\item 设 $A^2=A$（幂等矩阵），证明 $A$ 的特征值只能是 $0$ 或 $1$。
\item 证明：旋转矩阵 $\begin{pmatrix} 0 & -1 \\ 1 & 0 \end{pmatrix}$ 在实数
  范围内没有特征向量，并给出几何解释。
\item 证明：正交矩阵的乘积仍是正交矩阵，且正交矩阵的行列式为 $\pm1$。
\end{enumerate}
